\documentclass[12pt]{article}
\usepackage[utf8]{inputenc}
\usepackage[T1]{fontenc}
\usepackage{hyperref}
\title{Tugun Language Resources for K-12 Curriculum Development}
\author{Compiled by Claude Sonnet 5 (claude-sonnet-5) for Andr\'as Kornai}
\date{}
\begin{document}
\maketitle
\section{Introduction}
Tugun (tzn) is a Wetar Island language of Maluku, Indonesia, in the Austronesian family. Sources give about 1,200 speakers, and one gives about 1,800. Very little exists. No school materials were found.
\section{Primary Educational Resources}
None were found in the search results.
\section{Linguistic and Reference Materials}
The main description is Bryan Hinton's 2000 chapter in Charles Grimes's \emph{Spices from the East} (Pacific Linguistics), but no link to it was found. A Webonary site for Tugun has an introduction: \url{https://www.webonary.work/tugun/overview/introduction/} and an entry page \url{https://www.webonary.work/tugun/?p=9444}. The Wetarese article on Wikipedia places Tugun among the Wetarese varieties: \url{https://en.wikipedia.org/wiki/Wetarese_language}. Word lists and numerals are in the Austronesian Basic Vocabulary Database (\url{https://abvd.eva.mpg.de/austronesian/language.php?id=1639}) and the MPI numeral database (\url{https://lingweb.eva.mpg.de/channumerals/Tugun.htm}). Verbix has a language page: \url{https://docs.verbix.com/Languages/Tugun}.
\section{Religious and Community Resources}
None were found in the search results.
\section{Digital and Audio Resources}
Apart from the Webonary site, none were found.
\section{Curriculum Development Context}
There is no evidence of Tugun schooling. A curriculum would have to start from documentation work and community involvement, using Indonesian as the language of instruction in the meantime.
\end{document}